\bookchapter{Chapitre 29 : Santé numérique — Construire une relation saine avec votre smartphone}

\chapterintro{Avez-vous déjà vécu cela : vous vouliez juste regarder 5 minutes de courtes vidéos, et en levant les yeux, deux heures étaient passées ? Ou au lit le soir, plus vous regardez le smartphone, plus vous êtes réveillé, et finissez par vous endormir à 2 heures du matin ? Le smartphone a changé notre vie, mais il dévore aussi discrètement notre temps et notre santé. Ce chapitre ne vise pas à vous faire « arrêter le smartphone », mais à vous aider à reprendre le contrôle.}

\booksection{I. Regardez d'abord votre « temps d'écran » — les chiffres réveillent}

Avant de changer, il faut comprendre la situation. Le smartphone dispose d'une fonction que beaucoup n'ont jamais ouverte, appelée « Temps d'écran » (iPhone) ou « Bien-être numérique » (Android), qui vous indique combien de temps vous passez chaque jour sur le smartphone et sur quelles applications.

\begin{itemize}
  \item \textbf{iPhone}: Réglages > Temps d'écran > Voir toute l'activité. Vous verrez le nombre de déverrouillages quotidiens, le nombre de fois où vous prenez le smartphone, et la durée d'utilisation de chaque application.
  \item \textbf{Android}: Réglages > Bien-être numérique et contrôle parental. Huawei l'appelle « utilisation saine du smartphone », Xiaomi « gestion du temps d'écran ». Fonctions similaires, toutes affichent les données d'utilisation quotidiennes.
\end{itemize}

\textbf{80 \% des personnes qui consultent ces données pour la première fois sont choquées}. Un ami a découvert qu'il passait plus de 4 heures par jour sur Douyin, soit 28 heures par semaine — l'équivalent d'un emploi à temps partiel. Ce « choc numérique » est la première étape du changement.

\booksection{II. Définir des « limites d'utilisation » pour les applications — installez un limiteur de vitesse}

Une fois le problème identifié, des limites ciblées peuvent être définies. Pas besoin de supprimer l'application, fixez simplement un « quota quotidien » :

\begin{enumerate}
  \item \textbf{iPhone}: Réglages > Temps d'écran > Limites d'utilisation > Ajouter une limite. Sélectionnez l'application à limiter (par exemple Douyin), définissez 30 minutes par jour. À l'expiration, l'icône de l'application devient grise, et son ouverture nécessite un mot de passe supplémentaire.
  \item \textbf{Android}: Réglages > Bien-être numérique > Minuteur d'application. Sélectionnez l'application et définissez un plafond quotidien, à l'expiration l'icône devient grise.
  \item \textbf{Astuce}: Définissez un mot de passe de limite que vous n'utilisez pas souvent et difficile à mémoriser. S'il est trop facile à saisir, vous serez tenté de « prolonger » à chaque fois, et la limite n'aura plus aucun effet.
\end{enumerate}

\booksection{III. Bon usage du « mode Ne pas déranger » — ce n'est pas couper les notifications, c'est couper l'anxiété}

Le smartphone de beaucoup de gens sonne des dizaines de fois par jour, et chaque notification interrompt l'attention. Vous pensez « juste regarder la notification », mais en réalité, il faut 15-20 minutes pour se reconcentrer après chaque interruption. Le mode Ne pas déranger ne vous met pas « hors de portée », il vous aide à reprendre le contrôle de votre attention.

\begin{itemize}
  \item \textbf{Activation programmée}: Programmez le mode Ne pas déranger de 22h à 7h. Ainsi, le smartphone ne vous réveillera pas pour un SMS spam pendant la nuit.
  \item \textbf{Autoriser les contacts importants à passer}: Dans les réglages du mode Ne pas déranger, définissez « appels autorisés » sur « contacts favoris » ou « appels répétés » (un deuxième appel dans les 3 minutes sonnera normalement). Vous ne manquerez pas les appels urgents de la famille, mais ne serez pas dérangé par les appels commerciaux.
  \item \textbf{Mode concentration travail}: Le « mode Concentration » de l'iPhone permet de définir des scénarios personnalisés — par exemple, en mode « travail », seules les notifications des applications professionnelles sont autorisées ; en mode « lecture », toutes les notifications sont muettes.
\end{itemize}

\booksection{IV. Protéger ses yeux par l'habitude, pas par une application — la règle 20-20-20}

Regarder longtemps l'écran du smartphone provoque inévitablement une sécheresse et une fatigue oculaires. Le mode de protection des yeux (filtrage de la lumière bleue) aide un peu, mais ce qui est vraiment efficace, c'est de changer ses habitudes :

\begin{tcbinfo}{Règle 20-20-20}
Toutes les\textbf{20 minutes}de regard sur l'écran, déplacez votre regard vers\textbf{20 pieds (environ 6 mètres)}au loin, pendant\textbf{20 secondes}. Ce geste simple soulage efficacement la fatigue du muscle ciliaire.

\end{tcbinfo}

\begin{itemize}
  \item \textbf{Activer le mode protection des yeux}: Réglages > Affichage > Mode protection des yeux (ou mode Nuit), une fois activé, la température de couleur de l'écran devient plus chaude, réduisant la lumière bleue. Il est recommandé de le programmer du coucher au lever du soleil.
  \item \textbf{Ajuster la luminosité de l'écran}: Quand l'environnement est sombre, l'écran doit aussi être sombre. Beaucoup de gens regardent le smartphone dans le noir avec une luminosité trop élevée, c'est le comportement le plus nocif pour les yeux. Activez la « luminosité automatique » et assurez-vous qu'elle fonctionne.
  \item \textbf{Garder la distance}: Le smartphone doit être à au moins 30 cm des yeux. Beaucoup de gens le tiennent très près, ce qui maintient le muscle ciliaire en contraction permanente, accélérant la myopie.
\end{itemize}

\booksection{V. Guide de sauvetage cervical pour les «Accros au smartphone» (accros du smartphone penché)}

En regardant le smartphone la tête baissée, votre cou subit une pression 5 à 6 fois supérieure à la position debout normale. Une inclinaison prolongée entraîne une perte de la courbure cervicale, des ostéophytes, des vertiges et des engourdissements des mains. Voici quelques méthodes d'amélioration simples :

\begin{itemize}
  \item \textbf{Levez le smartphone}: En regardant le smartphone, essayez de le tenir à hauteur des yeux plutôt que de baisser la tête vers le smartphone posé sur vos genoux. Bien que les bras se fatiguent, votre cou vous remerciera.
  \item \textbf{Bougez le cou toutes les 30 minutes}: Tournez lentement la tête à gauche jusqu'au maximum, maintenez 3 secondes, revenez ; puis à droite, maintenez 3 secondes. Répétez 5 fois. Ensuite, faites l'« exercice du caractère Mi » — écrivez le caractère « Mi » (riz) dans les airs avec votre menton.
  \item \textbf{Debout contre un mur pendant 3 minutes}: Tête, omoplates, fesses et talons contre le mur, tenez debout 3 minutes par jour pour corriger efficacement la posture voûtée.
\end{itemize}

\booksection{VI. Sommeil et smartphone — l'habitude à changer dans la chambre}

La lumière bleue de l'écran du smartphone supprime la sécrétion de mélatonine, vous rendant « plus vous regardez, plus vous êtes réveillé ». Beaucoup de gens pensent « je m'endors en regardant », mais en réalité, le cerveau est trompé par la lumière bleue, et une fois le smartphone posé, il est difficile de s'endormir. Voici des conseils d'amélioration :

\begin{itemize}
  \item \textbf{Ne pas regarder le smartphone 1 heure avant le coucher}: C'est le conseil le plus efficace mais le plus difficile à suivre. Si vous n'y arrivez pas, activez au moins le « mode Nuit / protection des yeux » pour réduire au minimum la lumière bleue.
  \item \textbf{Placez le smartphone hors de portée}: Mettez le chargeur à l'autre bout de la chambre, pas sur la table de chevet. Ainsi, vous ne serez pas tenté de le prendre pour scroller, et quand l'alarme sonnera, vous devrez vous lever pour l'éteindre — d'une pierre deux coups.
  \item \textbf{Activez le « mode sommeil »}: Le mode de concentration « Sommeil » de l'iPhone peut verrouiller automatiquement l'écran pendant les heures de sommeil et simplifier l'écran de verrouillage ; les marques Android ont des fonctions similaires.
  \item \textbf{Utilisez un réveil au lieu du smartphone}: Achetez un réveil physique pour quelques dizaines de yuans, n'utilisez pas le smartphone comme réveil, vous découvrirez que le piège du « je regarde juste l'heure puis je scrolle » disparaît.
\end{itemize}

\booksection{VII. Anxiété informationnelle — apprenez la « cécité sélective »}

La quantité d'informations que nous recevons chaque jour équivaut à celle qu'une personne de l'Antiquité recevait en une vie entière. Mais le cerveau humain n'est pas conçu pour traiter autant d'informations ; la surcharge informationnelle entraîne anxiété, dispersion de l'attention et fatigue décisionnelle. La solution n'est pas « obtenir plus d'informations », mais « filtrer activement » :

\begin{itemize}
  \item \textbf{Se désabonner}: Ouvrez la liste de tous les comptes publics, créateurs et blogueurs que vous suivez, et posez-vous une question : « Au cours du mois écoulé, en ai-je tiré une valeur réelle ? » Si non, désabonnez-vous. Les canaux d'information doivent être qualitatifs, pas quantitatifs.
  \item \textbf{Désactiver les notifications non essentielles}: À l'exception de WeChat, téléphone, SMS et calendrier, désactivez toutes les notifications des autres applications. Les actualités n'ont pas besoin de notifications — vous les verrez naturellement quand vous irez les chercher.
  \item \textbf{Instaurez un « jour de jeûne informationnel »}: Choisissez un jour par semaine (par exemple le dimanche), sans actualités ni courtes vidéos sur le smartphone. Vous découvrirez que le ciel ne tombe pas, et que lundi votre état mental sera nettement meilleur.
\end{itemize}

\begin{tcbpractice}{Exercice : défi bien-être numérique d'une semaine}
\begin{enumerate}
  \item Aujourd'hui : ouvrez Temps d'écran / Bien-être numérique, consultez vos données de la semaine dernière.
  \item Demain : définissez un plafond quotidien pour l'application de divertissement la plus chronophage (30-60 minutes recommandées).
  \item Après-demain : programmez le mode Ne pas déranger (22h00-07h00), autorisez les appels des contacts favoris.
  \item Quatrième jour : ne regardez pas le smartphone 1 heure avant le coucher, placez le chargeur à l'autre bout de la chambre.
  \item Cinquième jour : désabonnez-vous de 5 comptes publics ou blogueurs que vous ne lisez plus.
  \item Sixième jour : essayez un « jeûne informationnel » complet — pas d'actualités, pas de courtes vidéos.
  \item Septième jour : consultez le rapport de temps d'écran de la semaine et comparez avec la semaine précédente.
\end{enumerate}

\end{tcbpractice}

\begin{tcbsummary}{Résumé du chapitre}
\begin{itemize}
  \item Consultez d'abord les données de temps d'écran, reconnaître la situation est la première étape du changement
  \item Fixez des limites aux applications de divertissement, programmez le mode Ne pas déranger, reprenez le contrôle de votre attention
  \item La règle 20-20-20 est l'habitude de protection des yeux la plus simple et la plus efficace
  \item En regardant le smartphone penché, le cou subit 5-6 fois le poids ; levez le smartphone et bougez le cou régulièrement
  \item Ne regardez pas le smartphone 1 heure avant le coucher, placez le chargeur hors de portée
  \item Le remède à l'anxiété informationnelle est le « filtrage actif », pas « obtenir plus »
\end{itemize}

\end{tcbsummary}
